% Curriculum Vitae -- Tae Soo Kim
% Build with XeLaTeX (e.g. `tectonic cv.tex`).
\documentclass[11pt,letterpaper]{article}

\usepackage[margin=1.5cm,bottom=52.02pt,includefoot,footskip=20.5pt]{geometry}
\usepackage{fontspec}
\usepackage[dvipsnames]{xcolor}
\usepackage{titlesec}
\usepackage{enumitem}
\usepackage{fancyhdr}
\usepackage{lastpage}
\usepackage[hidelinks]{hyperref}

% ---------------------------------------------------------------- colors
\colorlet{ink}{black!90}
\colorlet{muted}{black!60}
\colorlet{accent}{MidnightBlue}

% ---------------------------------------------------------------- page
\setlength{\parindent}{0pt}
\setlength{\parskip}{0pt}
\raggedbottom
\raggedright
\pagestyle{fancy}
\fancyhf{}
\renewcommand{\headrulewidth}{0pt}
\fancyfoot[R]{\color{ink}Page \thepage\ of \pageref*{LastPage}}
\AtBeginDocument{\color{ink}}

% ---------------------------------------------------------------- sections
\titleformat{\section}{\Large\scshape\color{accent}}{}{0pt}{}[{\titlerule}]
\titlespacing*{\section}{0pt}{16.5pt}{7.4pt}

% ---------------------------------------------------------------- lengths
\newlength{\entryindent}\setlength{\entryindent}{13.94pt}
\newlength{\itemindentx}\setlength{\itemindentx}{37.94pt}

% ---------------------------------------------------------------- entries
% \entry{Organization}{Location}{Role}{Dates}
\makeatletter
\newcommand{\entry}[4]{%
  \par\if@nobreak\vskip 2.3pt\else\vskip 8.3pt\fi
  {\leftskip=\entryindent\rightskip=1.87pt
   \textbf{#1}\hfill #2\par\nobreak\vskip 1.5pt
   {\small #3\hfill #4\par}}\nobreak}
\makeatother

% Dash list under an entry. Plain \item is a muted detail line;
% \head{Bold title}[trailing text] is a dark sub-heading line.
\newlist{cvitems}{itemize}{1}
\setlist[cvitems]{label={--}, leftmargin=\itemindentx, labelsep=5.45pt,
  topsep=4.8pt, partopsep=0pt, itemsep=4pt, parsep=0pt,
  before=\small\color{muted}, after=\vspace{-4.8pt}}
\NewDocumentCommand{\head}{m o}{%
  \item[{\normalsize\color{ink}\bfseries --}]%
  {\color{ink}\textbf{#1}\IfValueT{#2}{ #2}}}

% Bullet lists (teaching, skills, service)
% Same list in dark text, for education details
\newenvironment{cvitemsdark}
  {\begin{cvitems}[label={\bfseries --}, before=\small\color{ink}, topsep=4.1pt, after=\vspace{-3.8pt}]}
  {\end{cvitems}}

\newlist{cvbullets}{itemize}{1}
\setlist[cvbullets]{label=\textbullet, leftmargin=\entryindent, labelsep=5.4pt,
  topsep=2pt, partopsep=0pt, itemsep=4.4pt, parsep=0pt, before=\small}
% \dated{first line}{date}{second line}
\newcommand{\dated}[3]{\item #1\hfill #2\hspace*{3.3pt}\null\\[1.5pt]#3}

% ---------------------------------------------------------------- publications
\newcommand{\me}{\textbf{T.S.~KIM}}
% Keep each entry on one page.
\makeatletter
\newcommand{\pubsetup}{\clubpenalty=10000\@clubpenalty=10000\widowpenalty=10000}
\makeatother
\newlist{publications}{enumerate}{1}
\setlist[publications]{label={[\arabic*]}, ref={[\arabic*]}, align=right,
  leftmargin=27.9pt, labelwidth=17pt, labelsep=10.9pt,
  topsep=2.5pt, partopsep=0pt, itemsep=4.5pt, parsep=0pt, before={\frenchspacing\pubsetup}}
% \pub{key}{authors}{title}{venue}{year}
\newcommand{\pub}[5]{\item\label{pub:#1}#2, ``#3'', in \textit{#4}, #5.}
\newcommand{\pubref}[1]{\ref{pub:#1}}

% Bold label followed by a fixed gap (skills, languages)
\newcommand{\lbl}[1]{\textbf{#1}\hspace{9.4pt}}

\newcommand{\TM}{\textsuperscript{TM}}

\begin{document}

% ================================================================ header
\raisebox{-17.4pt}[0pt][0pt]{\Huge\color{accent}Tae Soo Kim}%
\hfill
\begin{minipage}[t]{0.45\textwidth}
  \raggedleft
  \href{mailto:taesoo.kim@lunit.io}{taesoo.kim@lunit.io}\\
  \href{https://taesoo-kim.github.io/}{https://taesoo-kim.github.io}\\
  \href{https://github.com/TaeSoo-Kim/}{github.com/TaeSoo-Kim}\\
  Last updated: November, 2026
\end{minipage}\hspace{22.7pt}

% ================================================================ objective
\section*{Objective}
Applied research scientist and team lead at Lunit with 10+ years across computer vision, medical AI and agentic systems.
I led a national consortium of 23 institutions and 70+ members (6 academic labs, 8 industry partners, 9 hospitals)
that built and open-sourced a medical foundation model, funded by the Korea Ministry of Science and ICT.
Before that, I led the team that shipped the AI models behind Lunit INSIGHT CXR, MMG and DBT.
I now lead the development of AI transformation (AX) products
that bring clinical-domain agents into everyday hospital workflows.

% ================================================================ education
\section*{Education}

\entry{Johns Hopkins University}{Baltimore, MD}
      {Ph.D. in Computer Science, Advisor: Dr. Gregory D. Hager}{2015--2021}
\begin{cvitemsdark}
  \item Thesis: ``Model-driven and Data-driven Methods for Recognizing Compositional Interactions from Videos''
\end{cvitemsdark}

\entry{Johns Hopkins University}{Baltimore, MD}
      {M.S. in Computer Science, Advisor: Dr. Russell Taylor and Dr. Austin Reiter}{2014--2015}
\begin{cvitemsdark}
  \item Thesis: ``3D Reconstruction System using a Flexible Endoscope and a Laser''
\end{cvitemsdark}

\entry{Johns Hopkins University}{Baltimore, MD}
      {B.S.E in Computer Science}{2010--2014}
\begin{cvitemsdark}
  \item Graduated with general honors
  \item Minor: Robotics
\end{cvitemsdark}

% ================================================================ professional
\section*{Professional Experience}

\entry{Lunit}{Seoul, Korea}
      {Staff Applied Research Scientist, Team Lead}{2021.06 - Current}
\begin{cvitems}
  \head{Staff Applied Research Scientist, AI Innovation Team Lead}[(2025.04 - Current)]
  \item The AI Innovation Team of the Technology Innovation Group is responsible for identifying and building Lunit's
        next generation of products by strategically securing necessary AI innovations.
  \item As a team lead, I contribute by leading a team of 10+ researchers and engineers to build multi-modal
        foundation models by securing funding opportunities by winning grants, scoping the product by meeting
        customers and overseeing the technical architecture of our services.
  \item I led the Lunit Consortium (6 academic labs, 8 industry partners, 9 hospitals) for the national
        `AI-Specialized Foundation Model' project, funded by the Korea Ministry of Science and ICT.
        \href{https://medigatenews.com/news/2040726640}{[See news article]}

  \head{Senior Applied Research Scientist, Model Team Lead}[(2022.09 - 2025.03)]
  \item The AI Model Team developed AI models for all product line ups in the Cancer Screening Department of Lunit
        including Lunit INSIGHT CXR\TM, Lunit INSIGHT MMG\TM\ and Lunit INSIGHT DBT\TM.
  \item As a team lead, I managed a team of 10+ researchers and set up scalable processes for data annotation,
        training, inference and evaluation. I initiated and drove a project to build multi-modal foundation models
        for chest X-rays, which resulted in new automated reporting products.
  \item We built evidence for AI algorithms for future Lunit products by executing forward looking research projects
        that directly improved our products and publishing scientific papers.
        (See publications: \pubref{spec-cxr}--\pubref{annotation-cost})

  \head{Senior Research Scientist, Chest X-ray Team Lead}[(2021.06 - 2022.09)]
  \item The Chest X-ray Team of the Cancer Screening Group is in charge of developing core deep learning algorithms
        used in the Lunit INSIGHT CXR\TM\ product line.
  \item As a research scientist and team lead, I contribute to the product by optimizing the vision model for lesion
        detection, adding new detectable lesions to the model, releasing AI engines and managing the team.
\end{cvitems}

\entry{Siemens Healthineers (Siemens Corporate Research)}{Princeton, NJ}
      {Graduate Research Intern: hosted by Dr. Shaohua Kevin Zhou}{Summer 2016}
\begin{cvitems}
  \head{Bone Removal Project}
  \item The aim of the Bone Removal project was to develop deep learning algorithms to remove voxels representing
        bones in a computed tomography angiography volume.
  \item As an integrated software for immediate deployment in products, the bone removal software was required to be
        highly precise, efficient and fast.
  \item As a graduate research intern, I developed a novel 3D convolutional neural network architecture with
        decomposed convolutions for bone removal.
  \item My approach increased accuracy and precision while decreasing memory usage and inference time. The
        developed bone removal model was successfully deployed to commercial products.
        (See patents: \pubref{patent-segmentation}, \pubref{patent-bone-removal}).
\end{cvitems}

\entry{Hyundai Heavy Industries, Medical Systems Department}{Seoul, Korea}
      {Undergraduate Research Intern}{Summer 2012}
\begin{cvitems}
  \head{Medical Robotics Project}
  \item The goal of the project was to build a minimally invasive surgical robot to remove malignant masses from breast
        cancer patients.
  \item As part of the computer vision team, I developed an algorithm for recognizing malignant/benign masses and
        cysts from breast ultrasound images. My main task involved data curation and implementation of state of the
        art machine learning algorithms.
\end{cvitems}

% ================================================================ academic
\section*{Academic Experience}

\entry{Johns Hopkins University}{Baltimore, MD}
      {Graduate Research Assistant}{2018-2021}
\begin{cvitems}
  \head{The Deep Intermodal Video Analytics (DIVA) project funded by IARPA}
  \item The DIVA program is a US government funded multi-institutional research program that aims to develop a
        robust automated activity detection system for multi-camera video surveillance applications.
  \item As a senior Ph.D student of the JHU DIVA team, I assumed a leadership role in the team's research,
        development and deployment of system deliverables to the program.
  \item I developed a zero-shot action recognition system that requires no labeled training examples of the test
        activities. (See publications: \pubref{daszl}, \pubref{safer})
  \item I developed a simulation software for synthesizing human activities from economic motion capture systems. We
        use the generated dataset to train activity recognition systems.
        (See publications: \pubref{synthetic-vehicles}, \pubref{synthesizing-attributes}, \pubref{unrealcv})
\end{cvitems}

\entry{Johns Hopkins University}{Baltimore, MD}
      {Graduate Research Assistant}{2016-2018}
\begin{cvitems}
  \head{Language of Surgery project funded by the Johns Hopkins Wilmer Eye Institute}
  \item In collaboration with the surgeons from the Johns Hopkins Wilmer Eye Institute, the project aimed to build
        intelligent tools to automatically analyze videos of cataract surgery.
  \item The goal is to build an intelligent system that provides actionable feedback to trainees by reviewing their
        surgical performance from video recordings.
  \item I showed that videos of cataract surgery can be annotated accurately at scale using crowd-sourced workers. (See
        publication: \pubref{crowdsourcing})
  \item As a machine learning and computer vision lead in the team, I developed deep learning algorithms to
        automatically recognize surgical phases from videos. (See publications: \pubref{cataract-phases})
  \item I developed a deep learning model that provides objective assessment of surgical technical skill from cataract
        surgery videos. (See publication: \pubref{capsulorhexis})
\end{cvitems}

\entry{Johns Hopkins University}{Baltimore, MD}
      {Masters Thesis Research}{2014-2015}
\begin{cvitems}
  \head{System for stereo reconstruction from monoscopic endoscope images}
  \item In collaboration with the department of head and neck surgery of the Johns Hopkins hospital, the project aimed
        to develop an economic solution for an endoscope capable of accurate 3D reconstruction of organs.
  \item I developed an active stereo vision system using a pattern generating laser inserted into a working channel of an
        endoscope.
  \item With a solution costing under a 100 dollars added to the endoscope, I developed a system capable of accurate
        3D reconstruction using monoscopic endoscope images. (See patent: \pubref{patent-endoscope} )
\end{cvitems}

% ================================================================ teaching
\section*{Teaching}
\begin{cvbullets}
  \dated{\textbf{Head Teaching Assistant} at Johns Hopkins University}{Fall 2015,2016}
        {\textbf{\textit{Computer vision}} \textit{(Dr. Austin Reiter)}}
  \dated{\textbf{Head Teaching Assistant} at Johns Hopkins University}{Spring 2016}
        {\textbf{\textit{Augmented Reality}} \textit{(Dr. Nasir Navab)}}
  \dated{\textbf{Head Teaching Assistant} at Johns Hopkins University}{Summer 2015}
        {\textbf{\textit{Intro. Programming}} \textit{(Dr. Sara More)}}
  \dated{\textbf{Head Teaching Assistant} at Johns Hopkins University}{Spring 2015}
        {\textbf{\textit{Data Structures}} \textit{(Dr. Sara More)}}
\end{cvbullets}

% ================================================================ skills / languages
\begin{minipage}[t]{0.49\textwidth}
  \section*{Skills}
  \begin{cvbullets}[itemsep=3.4pt]
    \raggedright
    \item \lbl{Machine Learning/Computer Vision:}Python,
          Pytorch, Keras, Tensorflow, Opencv, Matlab, Docker
    \item \lbl{Simulation:}UnrealEngine, UnrealCV, Axis-Neuron,
          ROS
    \item \lbl{Other:}Html, Php, Javascript, Java, SQL,
          Android-Studio, Confluence
  \end{cvbullets}
\end{minipage}%
\hfill
\begin{minipage}[t]{0.49\textwidth}
  \section*{Languages}
  \begin{cvbullets}
    \item \lbl{Korean:}Native
    \item \lbl{English:}Bilingual
  \end{cvbullets}
\end{minipage}

% ================================================================ service
\vspace{3.5pt}
\section*{Service}
\begin{cvbullets}
  \dated{\textbf{Organizer} of ICCV Multi-Modal Foundation Models for Cancer Detection and Prevention Workshop}{2025.10}
        {\textit{As part of the organization team, I contributed by preparing the proposal, recruiting event speakers and overseeing the
         program.}}
  \dated{\textbf{Organizer} of MICCAI Tutorial on AI for Medical Image Analysis in Practice}{2022.09}
        {\textit{As part of the organization team, I contributed by preparing the proposal, promoting the event and organizing the
         tutorial event at MICCAI.}}
  \dated{\textbf{President} of Korean Graduate Students Association at JHU}{2018--2019}
        {\textit{As an elected leader of the Korean graduate student body, I served to provide constructive support, opportunities as
         well as entertainment to members of the Korean community at JHU.}}
  \dated{\textbf{Reviewer}}{2015--current}
        {\textit{for CVPR, ICCV, ECCV, AAAI, MICCAI, IPCAI, TMI}}
\end{cvbullets}

% ================================================================ publications
\clearpage
\section*{Publications}
\begin{publications}
  \pub{deceptive-grounding}
      {C.~Caruzzo, D.~Yoo, and \me}
      {Deceptive grounding: Entity attribution failure in clinical retrieval-augmented generation}
      {NeurIPS}{2026}
  \pub{spec-cxr}
      {J.~Lee, J.~Cho, J.~Kim, L.~Dillard, T.~van~Sonsbeek, A.~Setio, H.~Lee, D.~Yoo, and \me}
      {SPEC-CXR: Advancing clinical safety through entity-level performance evaluation of chest X-ray report generation}
      {MICCAI}{2025}
  \pub{report-customization}
      {T.~Sonsbeek, A.~Setio, J.~Lee, J.~Cho, J.~Kim, H.~Lee, G.~Nam, L.~Dillard, and \me}
      {Beyond one size fits all: Customization of radiology report generation methods}
      {MICCAI-W}{2025}
  \pub{selectivekd}
      {L.~Dillard, H.~Lee, W.~Lee, \me, A.~Diba, and T.~Kooi}
      {SelectiveKD: A semi-supervised framework for cancer detection in DBT through knowledge distillation and pseudo-labeling}
      {MICCAI-W}{2024}
  \pub{rld}
      {J.~Song, \me, J.~Kim, G.~Nam, T.~Kooi, and J.~Choo}
      {Is user feedback always informative? Retrieval latent defending for semi-supervised domain adaptation without source data}
      {ECCV}{2024}
  \pub{prior-images}
      {H.~Lee, J.~Kim, E.~Park, M.~Kim, \me, and T.~Kooi}
      {Enhancing breast cancer risk prediction by incorporating prior images}
      {MICCAI}{2023}
  \pub{oooe}
      {G.~Nam, \me, S.~Lee, and T.~Kooi}
      {OOOE: Only-one-object-exists assumption to find very small objects in chest radiographs}
      {MICCAI-W}{2022}
  \pub{annotation-cost}
      {\me, G.~Jang, S.~Lee, and T.~Kooi}
      {Did you get what you paid for? Rethinking annotation cost of deep learning based computer aided detection in chest radiographs}
      {MICCAI}{2022}
  \pub{synthetic-vehicles}
      {\me, B.~Shim, M.~Peven, W.~Qiu, A.~Yuille, and G.~D.~Hager}
      {Learning from synthetic vehicles}
      {WACV-W Real World Surveillance}{2022}
  \pub{motion-guided-attention}
      {\me, J.~Jones, and G.~D.~Hager}
      {Motion guided attention fusion to recognize interactions from videos}
      {ICCV}{2021}
  \pub{daszl}
      {\me, J.~Jones, M.~Peven, Z.~Xiao, J.~Bai, Y.~Zhang, W.~Qiu, A.~Yuille, and G.~D.~Hager}
      {DASZL: Dynamic action signatures for zero-shot learning}
      {AAAI}{2021}
  \pub{synthesizing-attributes}
      {\me, M.~Peven, W.~Qiu, A.~Yuille, and G.~D.~Hager}
      {Synthesizing attributes with Unreal Engine for fine-grained activity analysis}
      {WACV-W}{2019}
  \pub{capsulorhexis}
      {\me, M.~O'Brien, S.~Zafar, A.~Malpani, G.~D.~Hager, S.~Sikder, and S.~Vedula}
      {Objective assessment of intraoperative technical skill in capsulorhexis using videos of cataract surgery}
      {IJCARS}{2019}
  \pub{safer}
      {\me, Y.~Zhang, Z.~Xiao, M.~Peven, W.~Qiu, J.~Bai, A.~Yuille, and G.~D.~Hager}
      {SAFER: Fine-grained activity detection by compositional hypothesis testing}
      {arXiv}{2019}
  \pub{cataract-phases}
      {\me*, F.~Yu*, G.~S.~Croso*, Z.~Song, F.~Parker, G.~D.~Hager, A.~Reiter, S.~Vedula, H.~Ali, and S.~Sikder}
      {Assessment of automated identification of phases in videos of cataract surgery using machine learning and deep learning techniques}
      {JAMA Network Open}{2019}
  \pub{crowdsourcing}
      {\me, A.~Malpani, A.~Reiter, G.~D.~Hager, S.~Sikder, and S.~Vedula}
      {Crowdsourcing annotation of surgical instruments in videos of cataract surgery}
      {MICCAI-W}{2018}
  \pub{unrealcv}
      {W.~Qiu, F.~Zhong, Y.~Zhang, S.~Qiao, Z.~Xiao, \me, Y.~Wang, and A.~Yuille}
      {UnrealCV: Virtual worlds for computer vision}
      {ACM-MM}{2017}
  \pub{tcn}
      {\me\ and A.~Reiter}
      {Interpretable 3D human action analysis with temporal convolutional networks}
      {CVPR-W}{2017}
  \pub{patent-segmentation}
      {S.~K.~Zhou, M.~Chen, H.~Ding, B.~Georgescu, M.~A.~Gulsun, \me, A.~P.~Kiraly, X.~Lu, J.-h.~Park, P.~Sharma, S.~Sun, D.~Xu, Z.~Xi, and Y.~Zheng}
      {Method and system for artificial intelligence based medical image segmentation}
      {US Patent Application Number 62/414,913}{2017}
  \pub{patent-bone-removal}
      {M.~Chen, \me, J.~Kretschmer, S.~Seifert, S.~K.~Zhou, M.~Schobinger, D.~Liu, Z.~Xu, S.~Grbic, and H.~Zhang}
      {Deep learning based bone removal in computed tomography angiography}
      {US Patent 10,852,907}{2019}
  \pub{patent-endoscope}
      {K.~C.~Olds, \me, R.~H.~Taylor, and A.~Reiter}
      {System for stereo reconstruction from monoscopic endoscope images}
      {US Patent 10,368,720}{2019}
\end{publications}

\end{document}
